% !PRINTED 336
\weberpage{336}

The coefficients of the cubic equation are known as soon as $m^3, mr^2, r^6$ are, or also when $m$ and $r^2$ are known.

If the equation (14) is then solved, $y_1,y_2,y_3$ are determined up to third roots of unity, which however must satisfy the relation $y_1y_2y_3=R$; and further than this these quantities cannot, by the nature of the matter, be determined.

If the curve $f$ is a real curve, then from the formulae (13) and (15) there follows further the theorem that, \gesperrt{when $m$ and $r$ are real, $D_1$ is positive, and consequently the roots of the equation (14) are all three real.}

If one expresses the discriminant $D_1$ in terms of the coefficients of the equation $P,Q,R$, then from the formula (15) there flows also the result that the square of the covariant $K$ can be expressed as an integral rational function of the covariants $f,\Delta,J$. The expression becomes rather complicated and shall not be pursued further here.

But in order to form the biquadratic equation on which $m$ and $r$ depend, it is necessary to enter more closely into the invariants of the third order.

\begin{center}
\S.\,91.

\gesperrt{The invariants of the curve of the third order and the biquadratic equation.}
\end{center}

All curves of the third order whose equation in the canonical form can be represented linearly by the same quantities $\Sigma y_1^3$, $y_1y_2y_3$ have the same points as points of inflection — for example, therefore, the curves $f=0$ and $\Delta=0$, and more generally all curves of the family
\weq{1}{F = \lambda f + 6\mu\Delta = 0,}
where $\lambda,\mu$ are arbitrary parameters. All curves of this family intersect in nine fixed points, and such a family is called a \gesperrt{pencil of curves}. The nine fixed points, which are the points of inflection for all curves of the pencil, are called the \gesperrt{fundamental points}. If we form from the form $F$ the \gesperrt{Hesse} determinant, this form too is expressed linearly by $\Sigma y_1^3$ and $y_1y_2y_3$, and can therefore by \S.\,90, (7) and (8) also be expressed linearly by $f$ and $\Delta$, i.e. we obtain a form of the same pencil.
